% Generated draft: requires mathematical and source review.
\documentclass[11pt]{article}
\usepackage[margin=25mm]{geometry}
\usepackage{fontspec}
\setmainfont{Latin Modern Roman}
\usepackage{amsmath,amssymb,mathtools,unicode-math}
\setmathfont{Latin Modern Math}
\usepackage{enumitem,needspace,xurl,hyperref}
\setlength{\parindent}{0pt}
\setlength{\parskip}{5pt}
\setlength{\emergencystretch}{3em}
\newcommand{\N}{\mathbb{N}}
\newcommand{\Z}{\mathbb{Z}}
\newcommand{\Q}{\mathbb{Q}}
\newcommand{\R}{\mathbb{R}}
\newcommand{\C}{\mathbb{C}}
\newcommand{\F}{\mathbb{F}}
\newcommand{\A}{\mathbb{A}}
\newcommand{\PP}{\mathbb{P}}
\newcommand{\E}{\mathbb{E}}
\newcommand{\eps}{\varepsilon}
\newcommand{\dd}{\,\mathrm{d}}
\begin{document}
\section*{UnsolvedMath ID: 1101102}
\textbf{Source problem code:} catalogue entry
\subsection{Definitions and mathematical statement}
Use the standard model
\[
 \operatorname{Sol}=\R^3,\qquad
 (x,y,z)(x',y',z')=(x+e^zx',\ y+e^{-z}y',\ z+z'),
\]
with left-invariant metric
\[
 \dd s^2=e^{-2z}\dd x^2+e^{2z}\dd y^2+\dd z^2.
\]
A self-map $F$ is a quasi-isometry if there are $K\geq1$, $C\geq0$ such that
\[
 K^{-1}d(u,v)-C\leq d(Fu,Fv)\leq Kd(u,v)+C
\]
for all $u,v$, and every point of Sol is at uniformly bounded distance from its image. Two such maps are at bounded distance if $\sup_u d(Fu,F'u)<\infty$. A map $f:\R\to\R$ is bi-Lipschitz if it is a bijection satisfying $L^{-1}|s-t|\leq|f(s)-f(t)|\leq L|s-t|$ for some $L\geq1$.

Kleiner's question asks for a classification of the quasi-isometries of this three-dimensional group, in particular a description of all self quasi-isometries up to bounded distance. The height coordinate is $h(x,y,z)=z$. The supplied assessment reports that every quasi-isometry is a bounded distance from a height-respecting one, with that approximation bound depending only on the quasi-isometry constants, and then describes bi-Lipschitz boundary-coordinate normal forms.

The assessment writes the latter as $(x,y,z)\mapsto(f(x),g(y),\pm z)$. Its negative-height branch needs clarification: in the standard metric, reversing height exchanges the expanding and contracting horizontal directions; the standard coordinate interchange is $(x,y,z)\mapsto(y,x,-z)$. The literal formula with unchanged coordinate roles must not be silently accepted as the negative-height normal form. Height translations and the distinction between an arbitrary finite approximation bound and a bound uniform in the quasi-isometry constants must also be kept explicit when a precise theorem is recovered from the cited papers.
\subsection{Short English statement}
Describe all quasi-isometries of Sol. The registry labels this solved, and the supplied assessment attributes height rigidity and a bi-Lipschitz product-type classification to Eskin--Fisher--Whyte, announced in 2007 with full proofs in 2012/2013. It says the original 2004 open item was thereby settled. Its abbreviated formula for the height-reversing case omits the coordinate-interchange convention, so the exact normal-form statement requires source review; no correction is silently promoted to the cited theorem here.
\subsection{Sources}
\begin{itemize}
\item UnsolvedMath permanent ID 1101102, source code catalogue entry, ``Questions in Geometric Group Theory --- Q 11.2.'' Registry curation and metadata: Ulam AI and UnsolvedMath Contributors, CC BY 4.0; updated 2026-08-06.
\item Bestvina, \emph{Questions in Geometric Group Theory} (source list dated 2004), Question 11.2, attributed to Kleiner, PDF page 19: \url{https://www.math.utah.edu/~bestvina/eprints/questions-updated.pdf}. Supplied manual transcription accessed 2026-07-29. Original metadata marks status and rights review NEEDS\_REVIEW; the later assessment claims to resolve the mathematical-status concern, not the rights review.
\item Eskin--Fisher--Whyte, 2007 announcement and full proofs in \emph{Annals of Mathematics} (2012/2013), as named in the supplied resolving assessment. The record provides no article titles, volumes, pages, or locators; none are invented.
\end{itemize}
\end{document}
